\documentclass[10pt,journal,twoside]{IEEEtran}

\usepackage[utf8]{inputenc}
\usepackage{amsmath,amsfonts,amssymb}
\usepackage{array}
\usepackage{textcomp}
\usepackage[table]{xcolor}
\usepackage{stfloats}
\usepackage{graphicx}
\usepackage{orcidlink}
\usepackage{cite}
\usepackage{multirow}
\usepackage{booktabs}
\usepackage{fancyhdr}
\hypersetup{
  colorlinks=true,
  linkcolor=blue,
  citecolor=blue,
  urlcolor=blue
}

\hyphenation{op-tical net-works semi-conduc-tor IEEE-Xplore audio-visual}
\raggedbottom

\fancypagestyle{tmmfirst}{%
  \fancyhf{}
  \fancyhead[L]{\scriptsize IEEE TRANSACTIONS ON MULTIMEDIA}
  \fancyhead[R]{\scriptsize\thepage}
  \renewcommand{\headrulewidth}{0pt}}
\fancypagestyle{tmmrunning}{%
  \fancyhf{}
  \fancyhead[LO]{\scriptsize WANG et al.: KA-GZSL: Knowledge-Guided Audio--Visual Generalized Zero-Shot Learning}
  \fancyhead[RE]{\scriptsize IEEE TRANSACTIONS ON MULTIMEDIA}
  \fancyhead[RO,LE]{\scriptsize\thepage}
  \renewcommand{\headrulewidth}{0pt}}
\pagestyle{tmmrunning}

\begin{document}

\title{KA-GZSL: Knowledge-Guided Audio–Visual Generalized Zero-Shot Learning}

% \author{Hu~Wang, Jing~Yang$^{*}$,~\IEEEmembership{Member,~IEEE}, Xiaoli~Ruan,Yuling~Chen, Xinru~Yi, Chengjiang~Li and Minyi~Guo,~\IEEEmembership{Fellow,~IEEE}
\author{
Hu~Wang\orcidlink{0009-0008-1596-5311},  % 替换成你的ORCID号
Jing~Yang$^{*}$\orcidlink{0000-0003-1915-9487},~\IEEEmembership{Member,~IEEE},
Xiaoli~Ruan\orcidlink{0000-0002-7623-3308},
Yuling~Chen\orcidlink{0000-0002-8674-8356},
Xinru~Yi\orcidlink{0009-0007-4224-4774},
Chengjiang~Li\orcidlink{0000-0002-1864-2023}
and
Minyi~Guo\orcidlink{0000-0003-0034-2302},~\IEEEmembership{Fellow,~IEEE}




\thanks{This work is supported by the Key Project of Guizhou Basic Research Program (QKHZD [2026] 047), Science and Technology Research Project of Guizhou Provincial Department of Education (No.QJJ[2025]014), National Natural Science Foundation of China under Grant (Nos. 62441608, 62166005), Scientific and Technological Innovation Platform Research Project of Guizhou Province (CXPTXM[2025]024),
Guizhou University Science and Technology Innovation Team ([2024] No.07),Guizhou University Basic Research Fund ([2024]08).(Corresponding author:Jing Yang.)\\
  \indent Hu Wang, Xiaoli Ruan, Yuling Chen and Minyi Guo are with the State Key Laboratory
 of Public Big Data, Guizhou University, Guiyang, 550025, China (e-mail: gs.wanghu24@gzu.edu.cn; xlruan@gzu.edu.cn; ylchen3@gzu.edu.cn; myguo@gzu.edu.cn).\\
 \indent Jing Yang is with the State Key Laboratory of Public Big Data,Guizhou University,Guiyang 550025,China,and also with the Department of Computer Science and Engineering,Shanghai Jiao Tong University,Shanghai 200240,China (e-mail:jyang23@gzu.edu.cn).\\
 \indent Xinru Yi is with the School of Mechanical Engineering, Guizhou University, Guiyang, Guizhou, 550025 China (e-mail:cme.xryi23@gzu.edu.cn)\\
  \indent Chengjiang Li is with the School of Management, Guizhou University, Guiyang 550025, China. (e-mail: cjli3@gzu.edu.cn).\\
  \indent  Minyi Guo is with the Department of Computer Science and Engineering, Shanghai Jiao Tong University, Shanghai 200240, China (e-mail: myguo@gzu.edu.cn).\\
  }
        % <-this % stops a space
%\thanks{This paper was produced by the IEEE Publication Technology Group. They are in Piscataway, NJ.}% <-this % stops a space
%\thanks{Manuscript received April 19, 2021; revised August 16, 2021.}
}

\maketitle
\thispagestyle{tmmfirst}
\begin{abstract}
Audio–visual generalized zero-shot learning aims to exploit the correlations among audio, visual, and semantic information to recognize and learn unseen classes. While AVCA and ClipClap-GZSL have made certain progress in terms of multimodal feature fusion, most existing methods rely on brief class labels for semantic modeling, making it difficult to adequately characterize fine-grained class semantics. Moreover, in real-world scenarios, these methods are vulnerable to weak audio–visual associations and modality noise. To address these issues, we propose knowledge-guided audio–visual generalized zero-shot learning (KA-GZSL), a method for aggregating external fine-grained knowledge and reducing distribution discrepancies across modalities. Specifically, we utilize large language models to perform fine-grained semantic expansion on class labels and combine them with the CLIP and CLAP text encoders to construct relational global action and relational global audio knowledge graphs, respectively, to enhance the semantic representation capability for unseen classes. Moreover, we design an audio–visual context encoder to fuse modality-specific and modality-shared information, thereby reducing the distribution discrepancy between different modalities. Furthermore, through the action knowledge graph and a cross-graph knowledge transfer mechanism, we structurally inject external prior knowledge into the audio–visual temporal features, explicitly modeling the deep correlations among audio, action, and visual features. The experimental results on three benchmark datasets—\text{VGGSound-GZSL}$^{\textit{cls}}$, \text{UCF-GZSL}$^{\textit{cls}}$, and \text{ActivityNet-GZSL}$^{\textit{cls}}$—demonstrate that the proposed method outperforms state-of-the-art methods across multiple metrics.
\end{abstract}

\begin{IEEEkeywords}
Audio–visual zero-shot learning, Video classification, Action knowledge graph, Audio–visual context encoder.
\end{IEEEkeywords}

\section{Introduction}
\IEEEPARstart{A}{udio--visual} content has become a major component of online
multimedia, and large-scale platforms continuously encounter new
actions and events that cannot be exhaustively annotated. Multimedia
applications such as video indexing, retrieval, recommendation, and
archive organization therefore require models that can recognize known
categories while generalizing to previously unseen ones. Audio--visual
generalized zero-shot learning (AVGZSL) addresses this need by
exploiting complementary audio and visual cues and associating
audio--visual representations with textual semantics, enabling the
joint recognition of seen and unseen classes. Representative methods,
including CJME~\cite{3}, AVGZSLNet~\cite{4}, AVCA~\cite{1}, and
ClipClap-GZSL~\cite{2}, have been developed to improve multimodal
representation and unseen-class recognition. Thus, AVGZSL constitutes
a multimedia representation problem involving heterogeneous signal
processing, cross-modal interaction, and semantic generalization.

Among these methods, AVCA and ClipClap-GZSL focus primarily on
improving feature fusion strategies. As shown in the representative
pipelines in Fig.~\ref{fig: Introduction}, existing approaches typically
employ Transformer architectures
(Fig.~\ref{fig: Introduction}(a)) or cross-attention mechanisms
(Fig.~\ref{fig: Introduction}(b)) to model temporal and cross-modal
interactions after feature extraction, and then align the resulting
representations with class-label embeddings in a joint space.
Nevertheless, two limitations remain: (i) text representations based
only on class labels provide insufficient fine-grained semantics,
leading to ambiguity among related action categories; and (ii) for the
weakly correlated audio--visual content prevalent in real-world videos,
dense cross-modal interactions without explicit knowledge guidance may
propagate irrelevant modality noise. These limitations are particularly
detrimental in open multimedia environments, where category
vocabularies evolve and audio--visual correspondence varies across
videos.

Motivated by these multimedia requirements, we reconsider the semantic
encoding process before feature fusion and argue that external
fine-grained knowledge is important for understanding emerging
concepts. To this end, we propose knowledge-guided audio--visual
generalized zero-shot learning (KA-GZSL). Specifically, an LLM-based
\textbf{K}nowledge-\textbf{A}ugmented \textbf{S}emantics (KAS)
component expands class labels into fine-grained descriptions,
compensating for the limited semantics provided by class names. We
further introduce an
\textbf{A}udio--\textbf{V}isual \textbf{C}ontext
\textbf{E}ncoder (AVCE) to integrate modality-specific evidence with
modality-shared temporal context in a common latent space. Finally, the
\textbf{A}ction \textbf{K}nowledge \textbf{G}raph (AKG) organizes
global class knowledge and sample-level audio--visual temporal features
as relational graphs and performs structured cross-graph knowledge
transfer. Together, these components provide a unified framework for
semantic enrichment, cross-modal interaction, and open-category
multimedia understanding.

To optimize the framework, we employ a weighted combination of joint
triplet, projection, and reconstruction losses. The joint objective
models the relationships among audio, visual, and action-semantic
features while preserving modality complementarity and enforcing
semantic consistency. Consequently, the model can transfer knowledge
from seen to unseen classes more effectively under complex multimedia
conditions.

\begin{figure}[!t]
    \centering
    \includegraphics[width=\columnwidth]{figure/Introduction.pdf}
    \caption{Comparison of representative AVGZSL pipelines. (a) Transformer-based methods emphasize intramodal representation, whereas (b) cross-attention methods emphasize audio--visual interaction; both commonly rely on short class-label embeddings. (c) KA-GZSL augments class semantics and transfers structured knowledge to modality-aware temporal representations.}
    \label{fig: Introduction}
\end{figure}


On the basis of the above design, the main contributions of this work
are summarized as follows:

\begin{itemize}
\item We utilize large language models to expand class semantics into
fine-grained, modality-oriented descriptions, alleviating the semantic
ambiguity caused by using class names alone and improving unseen-class
representation.

\item We propose an audio--visual context encoder that integrates
modality-specific cues with modality-shared temporal context, reducing
feature distribution discrepancies while preserving complementary
audio and visual information.

\item We introduce an action knowledge graph and a cross-graph
knowledge transfer mechanism to connect global class semantics with
sample-level audio--visual temporal graphs, providing structured
guidance under weak audio--visual correspondence.

\item Experiments on three AVGZSL benchmark datasets, namely,
VGGSound-GZSL$^{\textit{cls}}$, UCF-GZSL$^{\textit{cls}}$, and
ActivityNet-GZSL$^{\textit{cls}}$, demonstrate that KA-GZSL achieves
state-of-the-art performance and effectively supports open-category
multimedia recognition.
\end{itemize}


\section{Related Work}
\subsection{Audio–Visual Zero-Shot Learning}
Early AVGZSL methods learn coordinated audio--visual--semantic embeddings\cite{3,4}; AVCA\cite{1} and TCAF\cite{5} further model cross-modal and temporal interaction. Later work explores feature synthesis\cite{6}, hyperbolic embeddings\cite{7}, spiking temporal models\cite{8,43,44}, out-of-distribution detection\cite{13}, and tensor regression\cite{14}. Pretrained multimodal models provide stronger feature and text spaces: ClipClap-GZSL\cite{2} uses CLIP\cite{10} and CLAP\cite{11}, while EZ-GZSL\cite{12} simplifies training through optimized class embeddings. KDA\cite{9} enriches labels with LLM-generated descriptions, but the resulting semantics are still primarily aligned at the embedding level. In contrast, KA-GZSL uses enriched class semantics to construct fixed knowledge graphs and transfers their structure to sample-level temporal representations.

\subsection{Semantics-Guided Audio–Visual Learning}
Semantic guidance has improved multimedia representation in specialized settings. SGN transfers commentator semantics to figure-skating assessment\cite{19}, and language-guided models align domain knowledge with long-term audio--visual actions\cite{27}. Recent TMM work also uses foundation knowledge to separate sounding objects from interference\cite{liu2024bavs} and LLM-generated contextual views to distinguish visually similar video moments\cite{liu2025context}. These studies motivate semantic side information, but KA-GZSL addresses a different requirement: class-level knowledge must generalize from seen training videos to unseen action categories while preserving both modality-specific and shared evidence.

\subsection{Audio–Visual Learning and Graph Neural Networks}
Graphs provide explicit relational structure for multimedia learning, including visual scene graphs for source separation\cite{24}, motion-, sound-, and question-driven graphs for audio--visual question answering\cite{26}, and semantic graphs for event localization\cite{28}. Existing models generally build sample-dependent graphs for fully or weakly supervised tasks. KA-GZSL instead separates global class-knowledge graphs from per-video temporal graphs and performs one-way knowledge-to-temporal transfer, preventing individual samples from altering the class repositories used for unseen-category inference.

\begin{figure*}[!t]
    \centering
    \includegraphics[width=0.9\textwidth]{figure/KA-GZSL.pdf}
    \caption{Overview of the proposed KA-GZSL. Visual and audio features are first extracted separately from the input video. In KAS, an external LLM expands class labels into fine-grained descriptions, which are encoded by pretrained CLIP and CLAP text encoders. AVCE then integrates modality-specific features with shared temporal information in the latent space. Next, AKG injects global action semantics into audio–visual temporal features through GCN-based cross-graph knowledge transfer to model cross-modal constraints. Finally, the enhanced features are fused and projected into a common embedding space aligned with the textual semantics under joint loss optimization.}
    \label{fig:KA-GZSL}
\end{figure*}


\section{The Method}

\subsection{Problem definition}
At each evaluation stage, let $\mathcal{C}^{s}$ and $\mathcal{C}^{u}$ denote the disjoint classes with and without audio--visual training samples, respectively, where $\mathcal{C}^{s}\cap\mathcal{C}^{u}=\emptyset$. Their union is $\mathcal{C}=\mathcal{C}^{s}\cup\mathcal{C}^{u}$ and $N_c=|\mathcal{C}|$. The training set is $\mathcal{D}_{\mathrm{tr}}^{s}=\{(a_i^{s},v_i^{s},y_i^{s})\}_{i=1}^{N_{\mathrm{tr}}}$, with $y_i^{s}\in\mathcal{C}^{s}$; no audio--visual sample from the held-out $\mathcal{C}^{u}$ at that stage is used for training. Standard ZSL evaluates $\mathcal{C}^{u}$, whereas GZSL evaluates $\mathcal{C}^{s}\cup\mathcal{C}^{u}$. Class-level textual side information is available for all classes. In the two-stage benchmark protocol, validation classes serve as pseudo-unseen classes in stage one and join the seen training set in stage two; the final test-unseen classes remain held out.



\subsection{Overall framework}
\textbf{Overall architecture.} In conventional methods, the text branch typically relies only on class labels to generate semantic embeddings, which are used only for the optimization of the final loss function. In contrast, our method leverages rich textual descriptions of the class labels and combines them with graph convolutional networks to construct a knowledge graph, thereby guiding and promoting the fusion of the audio and visual modalities. As shown in Fig.~\ref{fig:KA-GZSL}, the proposed framework consists of three main components: the KAS module, the AKG module, and the AVCE module. First, during the feature extraction stage, the KAS module provides fine-grained descriptions of each class to enrich the category semantics, thereby alleviating the issue of the limited expressive capability of the class name embeddings in previous methods. Second, the AVCE module is designed to fuse modality-specific features and modality-shared features, addressing the modality interference and semantic loss caused by direct concatenation and producing more discriminative cross-modal representations. In the relational graph construction stage, the AKG module employs graph convolutional networks to fuse the action knowledge and audio–visual temporal features in a structured manner, thereby mitigating the weak correlations between the visual and audio modalities and the poor coupling introduced by naively concatenating features with knowledge.

\subsection{Knowledge-Augmented Semantics}
In audio–visual generalized zero-shot learning, the model must transfer knowledge learned from seen classes to unseen classes. However, existing methods are limited by coarse-grained semantic representations and ambiguous action concepts. For example, for semantically similar unseen classes such as playing basketball and slam dunk, it is difficult to distinguish the differences among various actions using only the class labels. Inspired by the way in which humans leverage knowledge bases to improve our understanding of novel visual concepts, we propose the KAS module, which uses a large language model to transform abstract class labels into concrete and detailed action semantic descriptions. Specifically, by designing task-specific prompt templates, we guide the LLM to generate detailed descriptions containing visual, motion, and auditory characteristics, thereby enhancing the model’s ability to understand and discriminate unseen classes.



To address the issue that different modalities emphasize various aspects of semantics (for instance, in playing basketball, the visual modality focuses more on the dribbling motion, whereas the audio modality emphasizes the sound of the ball striking the floor), we construct a global visual action semantic repository and a global audio action semantic repository. For each class $y_n\in\mathcal{C}$, we instantiate a visual prompt $p_n^v$ and an audio prompt $p_n^a$ from the generic templates “a photo of a person doing {category}” and “a person doing {category} can be heard,” respectively. These modality-specific prompts introduce visual and auditory action knowledge for all $N_c$ classes.


Each class label is subsequently fed into the LLM to generate a detailed natural language description. Let $m\in\{v,a\}$ denote the visual or audio modality. The generated description is combined with the corresponding modality-specific prompt to characterize the visual appearance, motion pattern, and auditory properties of that class:
\begin{equation}
D_n^m=\operatorname{ConcatText}\left(p_n^m,\operatorname{LLM}(y_n)\right)
\end{equation}
where $n=1,\ldots,N_c$ and $m\in\{v,a\}$. The operator $\operatorname{LLM}(\cdot)$ generates a detailed class description, and $\operatorname{ConcatText}(\cdot,\cdot)$ denotes textual concatenation. For example, the LLM can be prompted with “Please provide a detailed explanation of the body movements involved in playing basketball.” Here, $D_n^m$ denotes the modality-oriented description of class $y_n$.


For the textual embeddings of the class labels, we use modality-specific text encoders:
\begin{equation}
f_n^{\mathrm{txt},m}=E_{\mathrm{txt}}^m\left(D_n^m\right),\quad n=1,\ldots,N_c,\quad m\in\{v,a\}
\end{equation}
where $E_{\mathrm{txt}}^v$ and $E_{\mathrm{txt}}^a$ are the frozen CLIP and CLAP text encoders, respectively, and $f_n^{\mathrm{txt},m}$ is the modality-oriented semantic representation of class $y_n$.


To capture local audio--visual correspondence, each video is divided into $T$ nonoverlapping segments of 32 consecutive frames~\cite{32}. Let $x_t^m$ denote the $t$-th segment of modality $m\in\{v,a\}$, where $v$ and $a$ represent the visual and audio modalities, respectively. Segment-level features are extracted by modality-specific pretrained encoders:
\begin{equation}
F_t^m=E_m\left(x_t^m\right),\quad t=1,\ldots,T,\quad m\in\{v,a\}
\end{equation}
where $E_v$ and $E_a$ denote the CLIP image encoder and the CLAP audio encoder, respectively, and $F_t^m$ is the feature of the $t$-th segment. Audio segments are preprocessed before encoding.

\subsection{Audio–Visual Context Encoding}
In multimodal learning, the effective use of the rich and diverse information from different modalities remains a key issue. Although modality-specific features can provide clean and diverse modality-aware information, under language supervision, these features may still lead to shallow unimodal learning. More specifically, our goal is to learn the latent associations among audio, action, and visual features rather than independently modeling isolated action–audio and action–visual relationships. To this end, we design an audio–visual context encoder to fuse modality-specific information with modality-shared information.


For modality-specific information, frozen pretrained encoders are adopted to extract the raw modality features, thereby preserving the fundamental representation capability of each modality. Specifically, the ViT-B/32 encoder of CLIP is used to extract the visual-specific features $\left\{ F_t^v \right\}_{t=1}^T$, while the audio encoder of CLAP is employed to extract the audio-specific features $\left\{ F_t^a \right\}_{t=1}^T$.


We subsequently employ trainable token projection networks to map the modality-specific features to a shared $d$-dimensional latent space:
\begin{equation}
P_t^m=\operatorname{Proj}_m\left(F_t^m\right),\quad t=1,\ldots,T,\quad m\in\{v,a\}
\end{equation}
where $\operatorname{Proj}_m(\cdot)$ is a two-layer multilayer perceptron for modality $m$.


For the modality-shared information, the projected features are fed into a shared encoder $\text{Encoder}_s$:
\begin{equation}
\left\{Z_t^m\right\}_{t=1}^{T}=\text{Encoder}_s\left(\left\{P_t^m\right\}_{t=1}^{T}\right),\quad m\in\{v,a\}
\end{equation}
where $\text{Encoder}_s(\cdot)$ denotes the shared transformer encoder. The shared transformer encoder\cite{33} captures long-range and modality-invariant temporal context across segments.


Finally, the modality-specific and modality-shared features are combined via elementwise summation:
\begin{equation}
f_t^m=P_t^m\oplus Z_t^m,\quad t=1,\ldots,T,\quad m\in\{v,a\}
\end{equation}
where $\oplus$ denotes elementwise addition and $f_t^m$ is the fused feature of the $t$-th segment for modality $m$.


\subsection{Action Knowledge Graph}
To address the limitations of previous methods, in which audio–visual features and class knowledge are fused merely through simple concatenation, resulting in weak semantic associations, we introduce the action knowledge graph to fuse the knowledge representations and audio–visual temporal features in a structured manner. On the basis of graph convolutional networks (GCNs), this module produces two independent knowledge–temporal interaction graphs for the visual and audio modalities. Specifically, the visual branch captures the temporal continuity of the visual features and their associations with the global visual semantics, whereas the audio branch captures the rhythmic variation patterns of the audio features and their associations with the global audio semantics. Through cross-graph message passing, the AKG enables deep integration between the knowledge representations and features, thereby improving the model’s semantic understanding and discriminative ability for unseen classes. This process involves two key steps: graph construction and cross-graph knowledge transfer.

\subsubsection{Graph Construction}
Let $m\in\{v,a\}$ index the visual and audio modalities. The modality-specific text projection $\operatorname{Proj}_{\mathrm{txt},m}$ maps each class semantic feature to the same $d$-dimensional space as the AVCE output. The initial class-knowledge node matrix and temporal node matrix are defined as
\begin{equation}
K_{\mathrm{act}}^{m,(0)} = \operatorname{Stack}_{n=1}^{N_c}\left(\operatorname{Proj}_{\mathrm{txt},m}\left(f_n^{\mathrm{txt},m}\right)\right)\in\mathbb{R}^{N_c\times d}
\end{equation}
\begin{equation}
H_m^{(0)} = \operatorname{Stack}_{t=1}^{T}\left(f_t^m\right)\in\mathbb{R}^{T\times d}
\end{equation}
where $\operatorname{Stack}$ arranges its vector arguments row-wise. Thus, $K_{\mathrm{act}}^{v,(0)}$ and $K_{\mathrm{act}}^{a,(0)}$ form the node matrices of the relational global action knowledge graph $G_{\mathrm{act}}^v$ and the relational global audio knowledge graph $G_{\mathrm{act}}^a$, respectively, whereas $H_v^{(0)}$ and $H_a^{(0)}$ form the node matrices of the visual and audio temporal graphs $G_v$ and $G_a$, respectively.

For the two knowledge graphs, the semantic relation between class nodes $n$ and $n'$ is measured by cosine similarity. For the two temporal graphs, the relation between segments $t$ and $t'$ is computed in the same manner:
\begin{equation}
\rho_{\mathrm{act},nn'}^m
=\frac{\left(f_n^{\mathrm{txt},m}\right)^{\top}f_{n'}^{\mathrm{txt},m}}{\left\|f_n^{\mathrm{txt},m}\right\|_2\left\|f_{n'}^{\mathrm{txt},m}\right\|_2}
\end{equation}
\begin{equation}
\rho_{tt'}^m
=\frac{\left(f_t^m\right)^{\top}f_{t'}^m}{\left\|f_t^m\right\|_2\left\|f_{t'}^m\right\|_2}
\end{equation}
Each class node is connected to its $k_{\mathrm{cls}}$ most similar class nodes. Each temporal node is connected both to its immediate temporal neighbors and to its $k_{\mathrm{tmp}}$ most similar nonself temporal nodes:
\begin{equation}
\mathcal{E}_{\mathrm{act}}^m=
\left\{(n,n')\mid n'\in\operatorname{KNN}_{k_{\mathrm{cls}}}
\left(n;\rho_{\mathrm{act}}^m\right)\right\}
\end{equation}
\begin{equation}
\begin{aligned}
\mathcal{E}_m={}&\left\{(t,t')\mid 0<|t-t'|\le 1\right\}\\
&\cup\left\{(t,t')\mid t'\in\operatorname{KNN}_{k_{\mathrm{tmp}}}
\left(t;\rho^m\right)\right\}
\end{aligned}
\end{equation}
Here, $\operatorname{KNN}_k(i;\rho)$ returns the indices of the $k$ nodes $q\ne i$ with the largest similarities $\rho_{iq}$. We use $k_{\mathrm{cls}}=5$ and $k_{\mathrm{tmp}}=\min(3,T-1)$. This hybrid temporal topology preserves local continuity while also capturing nonlocal segments with similar visual content or acoustic patterns. No additional similarity threshold is used; the top-$k$ rules provide a consistent sparsity level across datasets.

The raw edge-weight matrices $B_{\mathrm{act}}^m\in\mathbb{R}^{N_c\times N_c}$ and $B_m\in\mathbb{R}^{T\times T}$ are obtained by shifting the retained cosine similarities to the nonnegative interval $[0,1]$:
\begin{equation}
\left(B_{\mathrm{act}}^m\right)_{nn'}=
\begin{cases}
\dfrac{1+\rho_{\mathrm{act},nn'}^m}{2} & (n,n')\in\mathcal{E}_{\mathrm{act}}^m\\
0 & \text{otherwise}
\end{cases}
\end{equation}
\begin{equation}
\left(B_m\right)_{tt'}=
\begin{cases}
\dfrac{1+\rho_{tt'}^m}{2} & (t,t')\in\mathcal{E}_m\\
0 & \text{otherwise}
\end{cases}
\end{equation}
Because top-$k$ neighborhoods can be asymmetric, the raw matrices are symmetrized before self-loops are added. We define the symmetrized self-loop matrices as $R_{\mathrm{act}}^m=\max(B_{\mathrm{act}}^m,(B_{\mathrm{act}}^m)^{\top})+I_{N_c}$ and $R_{\mathrm{tmp}}^m=\max(B_m,B_m^{\top})+I_T$. The normalized knowledge-graph and temporal-graph adjacency matrices are computed separately:
\begin{equation}
\widehat A_{\mathrm{act}}^m = \left[\operatorname{diag}\left(R_{\mathrm{act}}^m\mathbf{1}_{N_c}\right)\right]^{-1/2}
R_{\mathrm{act}}^m
\left[\operatorname{diag}\left(R_{\mathrm{act}}^m\mathbf{1}_{N_c}\right)\right]^{-1/2}
\end{equation}
\begin{equation}
\widehat A_{\mathrm{tmp}}^m = \left[\operatorname{diag}\left(R_{\mathrm{tmp}}^m\mathbf{1}_{T}\right)\right]^{-1/2}
R_{\mathrm{tmp}}^m
\left[\operatorname{diag}\left(R_{\mathrm{tmp}}^m\mathbf{1}_{T}\right)\right]^{-1/2}
\end{equation}
Here, $\max(\cdot,\cdot)$ is applied elementwise, $I_{N_c}$ and $I_T$ add self-loops, and $\operatorname{diag}(R_{\mathrm{act}}^m\mathbf{1}_{N_c})$ and $\operatorname{diag}(R_{\mathrm{tmp}}^m\mathbf{1}_T)$ are the corresponding degree matrices. The knowledge-graph topology is constructed once from the frozen class text embeddings and remains fixed. The temporal-graph topology is deterministically recomputed for each input from the AVCE features and then shared by all AKG layers. These adjacency matrices are not free trainable parameters; $\operatorname{Proj}_{\mathrm{txt},m}$ and the graph-convolution weights are learned during training. All $N_c$ class-level text nodes, including unseen-class nodes, participate in the knowledge graphs, but no audio--visual sample from an unseen class is used to construct either the knowledge graphs or the temporal graphs.

\subsubsection{Cross-Graph Knowledge Transfer}
The purpose of the AKG is to use the global, class-level semantic repositories to guide sample-level temporal representations without allowing individual samples to alter the global class prototypes. We therefore employ two parallel one-way transfer streams: $G_{\mathrm{act}}^v\rightarrow G_v$ and $G_{\mathrm{act}}^a\rightarrow G_a$. For the $l$-th graph-transfer layer, where $l=0,\ldots,L_g-1$, intragraph propagation is first performed in both the knowledge graph and the corresponding temporal graph:
\begin{equation}
K_{\mathrm{act}}^{m,(l+1)}=
\sigma\left(
\widehat A_{\mathrm{act}}^mK_{\mathrm{act}}^{m,(l)}W_{\mathrm{act},m}^{(l)}
\right)
\end{equation}
\begin{equation}
\widetilde H_m^{(l+1)}=
\sigma\left(
\widehat A_{\mathrm{tmp}}^mH_m^{(l)}W_{\mathrm{tmp},m}^{(l)}
\right)
\end{equation}
where $W_{\mathrm{act},m}^{(l)}$ and $W_{\mathrm{tmp},m}^{(l)}$ are learnable graph-convolution weights, $L_g$ is the number of graph-transfer layers, and $\sigma(\cdot)$ is the ReLU activation. The intermediate representation $\widetilde H_m^{(l+1)}$ contains intragraph temporal context before class knowledge is injected.

The updated temporal nodes serve as queries, and the updated class-knowledge nodes serve as keys and values. Their row-wise knowledge-to-temporal mapping is
\begin{equation}
\Lambda_{\mathrm{act}\rightarrow m}^{(l+1)}=
\operatorname{Softmax}_{N_c}\left(
\frac{\widetilde H_m^{(l+1)}\left(K_{\mathrm{act}}^{m,(l+1)}\right)^{\top}}{\sqrt{d}}
\right)
\in\mathbb{R}^{T\times N_c}
\end{equation}
where $\operatorname{Softmax}_{N_c}$ normalizes every temporal-node row over the $N_c$ class-knowledge nodes. The current layer then injects the updated knowledge representations into the updated temporal representations:
\begin{equation}
H_m^{(l+1)}=
\sigma\left(
\widetilde H_m^{(l+1)}+
\Lambda_{\mathrm{act}\rightarrow m}^{(l+1)}
K_{\mathrm{act}}^{m,(l+1)}W_{\mathrm{cross},m}^{(l)}
\right)
\end{equation}
where $W_{\mathrm{cross},m}^{(l)}$ is a learnable knowledge projection. The first term preserves the intragraph temporal context, whereas the second term aggregates class-level semantic knowledge according to the relevance of each class node to each temporal segment. No reverse temporal-to-knowledge update is performed; consequently, the class repositories remain global and sample-independent while still providing modality-specific guidance to the visual and audio streams.


\subsection{Projection to the Joint Embedding Space}
After $L_g$ graph-transfer layers, the visual and audio temporal node matrices are $H_v^{(L_g)}\in\mathbb{R}^{T\times d}$ and $H_a^{(L_g)}\in\mathbb{R}^{T\times d}$, respectively. We denote their $t$-th rows by $h_{v,t}^{(L_g)}$ and $h_{a,t}^{(L_g)}$. These node representations contain both temporal context and class-level semantic knowledge. Global average pooling over the $T$ temporal nodes yields
\begin{equation}
\bar{v}=\frac{1}{T}\sum_{t=1}^{T}h_{v,t}^{(L_g)},\qquad
\bar{a}=\frac{1}{T}\sum_{t=1}^{T}h_{a,t}^{(L_g)}
\end{equation}
Here, $\bar{v}$ and $\bar{a}$ are the aggregated global visual and audio vectors, respectively; $T$ is the number of temporal segments and is distinct from the class-node count $N_c$.


Since the primary objective of AVGZSL is to predict the textual label category of the input sample, we employ a projection layer and a reconstruction layer to map the joint audio–visual embeddings to a unified space and reconstruct the projected features to preserve the original information. This design ensures the comparability of the features from different modalities. To fully exploit multimodal complementarity, the visual and audio vectors are concatenated and projected to a common space with the same dimensionality as the semantic embeddings via a fully connected layer as follows:
\begin{equation}
O_{\text{av}} = W_{\text{proj}} \cdot \left[ \bar{v} \parallel \bar{a} \right] + b_{\text{proj}}
\end{equation}
where $\left[ \cdot \parallel \cdot \right]$ denotes the feature concatenation operation, $W_{\text{proj}}$ and $b_{\text{proj}}$ are the learnable projection weights and bias, respectively, and $O_{\text{av}}$ is the final audio–visual hybrid representation used for classification. Similarly, the class semantic embeddings are projected to obtain $\{O_w^n\}_{n=1}^{N_c}$. During inference, the predicted class is determined by the projected class semantic embedding nearest to the projected audio--visual representation:
\begin{equation}
\widehat n=\operatorname*{arg\,min}_{1\le n\le N_c}\left\|O_w^n-O_{\text{av}}\right\|_2
\end{equation}
where $O_w^n$ is the projected semantic representation of class $y_n$, $\widehat n$ is the predicted class index, and $\widehat y=y_{\widehat n}$ is the corresponding class label.


\subsection{Loss function}
The model is trained by optimizing the overall loss $L$. It comprises the joint triplet loss $L_n$, projection loss $L_p$, and reconstruction loss $L_c$.


Joint Triplet Loss. The joint triplet loss $L_n$ plays a crucial role in clustering the final audio–visual embeddings, ensuring both compactness within classes and separability across classes. It is formally defined as follows:
\begin{equation}
L_n = \left[ \gamma + O_{\text{av}}^+ - O_w^+ \right]_+ + \left[ \gamma + O_{\text{av}}^- - O_w^+ \right]_+
\end{equation}
where $\gamma$ is the minimum margin between the mismatched negative pairs from different modalities and the correctly matched audio–visual embedding. Here, $O_w$ denotes the textual embedding, $O_{\text{av}}^+$ and $O_{\text{av}}^-$ denote the positive and negative samples, respectively, and $\left[ x \right]_+ \equiv \max(x,0)$.


Projection Loss. The projection loss $L_p$ minimizes the distance between each projected audio--visual embedding and the semantic embedding of its ground-truth class:
\begin{equation}
L_p=\frac{1}{N_{\mathrm{mb}}}\sum_{i=1}^{N_{\mathrm{mb}}}
\left\|O_{\mathrm{av},i}-O_w^{y_i}\right\|_2^2
\end{equation}
where $N_{\mathrm{mb}}$ is the minibatch size, $i$ indexes samples within the minibatch, and $O_w^{y_i}$ is the projected semantic representation of the ground-truth class $y_i$.


Reconstruction Loss. The reconstruction loss $L_c$ aims to preserve the original data distribution when the features are embedded in the unified space. The architecture of the reconstruction layer is similar to that of the projection layer, and the loss is formulated as follows:
\begin{equation}
L_c=\frac{1}{N_{\mathrm{mb}}}\sum_{i=1}^{N_{\mathrm{mb}}}
\left\|O_{\mathrm{av},i}^{\mathrm{rec}}-O_{w,y_i}^{\mathrm{rec}}\right\|_2^2
\end{equation}
where $O_{\mathrm{av},i}^{\mathrm{rec}}$ and $O_{w,y_i}^{\mathrm{rec}}$ are the reconstructed audio--visual and class-semantic representations associated with sample $i$, respectively.


The overall loss function is formulated as follows:
\begin{equation}
\label{eq:total_loss}
L=\alpha L_n+\beta L_p+\varphi L_c
\end{equation}
where $\alpha$, $\beta$, and $\varphi$ are the weighting coefficients of the corresponding loss terms. The three terms jointly constrain discriminative separation, semantic projection, and reconstruction.


\section{Experiment}



\subsection{Experimental setup}
\subsubsection{Dataset}
We evaluate KA-GZSL on three AVGZSL benchmarks. \text{VGGSound-GZSL}$^{\textit{cls}}$, derived from VGGSound\cite{34}, covers diverse sound-producing events; \text{UCF-GZSL}$^{\textit{cls}}$, derived from UCF101\cite{35}, focuses on daily, sports, and entertainment actions; and \text{ActivityNet-GZSL}$^{\textit{cls}}$, derived from ActivityNet\cite{36}, contains temporally annotated daily activities. Detailed dataset descriptions and complete class partitions under the two-stage protocol are provided in the supplementary material.


\subsubsection{Implementation details}
The model is implemented in PyTorch and trained on an NVIDIA GeForce RTX 3090Ti GPU using Adam\cite{37} with a batch size of 64. The loss weights in Eq. \eqref{eq:total_loss} are $\alpha=0.5$, $\beta=0.7$, and $\varphi=0.5$. We follow the two-stage protocol\cite{1,5,7}: the first stage uses $\text{Val}_\text{S}$ and $\text{Val}_\text{U}$ to select the calibrated-stacking parameter and epoch count, and the second stage retrains on $\{\text{train} \cup \text{Val}_\text{S} \cup \text{Val}_\text{U}\}$ before final testing. Other settings follow ClipClap-GZSL\cite{2}. The complete hardware configuration, optimizer parameters, learning-rate schedules, and dataset-specific epoch counts are provided in the supplementary material.


\subsubsection{Evaluation metrics}
We evaluate the performance of the model on seen and unseen classes using the S, U, HM, and ZSL metrics. Specifically, we report the average accuracy for seen classes (S) and unseen classes (U). In the context of GZSL, the model tends to exhibit an inherent bias toward seen categories, which often results in higher accuracy for seen classes. To balance S and U, we utilize the harmonic mean (HM) of S and U as the evaluation metric for the GZSL task. For ZSL, only classes from the unseen test class subset are selected as prediction results. Thus, in this task, only the subset of test samples from unseen classes is used in the performance evaluation.

\newlength{\oldaboverulesep}
\newlength{\oldbelowrulesep}
\begin{table*}[!t]
  \centering
  \caption{Performance of our model compared to state-of-the-art methods for AVGZSL on the \text{VGGSound-GZSL}$^{\textit{cls}}$, \text{UCF-GZSL}$^{\textit{cls}}$, and \text{ActivityNet-GZSL}$^{\textit{cls}}$ datasets. The overall top results are indicated in \textbf{bold} and the second-best results are \underline{underline}.}
  \label{tab:sota}
  
  \begin{tabular}{l cccc c cccc c cccc}
    \toprule[2pt]
    \multirow{2}{*}{\textbf{Methods}} 
    & \multicolumn{4}{c}{\text{VGGSound-GZSL}$^{\textit{cls}}$} & & \multicolumn{4}{c}{\text{UCF-GZSL}$^{\textit{cls}}$} & & \multicolumn{4}{c}{\text{ActivityNet-GZSL}$^{\textit{cls}}$} \\
    \cmidrule(lr){2-5} \cmidrule(lr){7-10} \cmidrule(lr){12-15}
    & S & U & HM & ZSL & & S & U & HM & ZSL & & S & U & HM & ZSL \\
    \midrule[1pt]
    % 第一列：方法名正常，会议名缩小；数值列：加粗最高，下划线次高
    DeViSE \cite{38} \,{\scriptsize \textbf{NeurIPS'13}} & \underline{36.22} & 1.07 & 2.08 & 5.59 & & 55.59 & 14.94 & 23.56 & 16.09 & & 3.45 & 8.53 & 4.91 & 8.53 \\
    ALE \cite{39} \,{\scriptsize \textbf{T-PAMP'15}} & 0.28 & 5.48 & 0.53 & 5.48 & & 57.59 & 14.89 & 23.66 & 16.32 & & 2.63 & 7.87 & 3.94 & 7.90 \\
    SJE \cite{40} \,{\scriptsize \textbf{CVPR'20}} & \textbf{48.33} & 1.10 & 2.15 & 4.06 & & 63.10 & 16.77 & 26.50 & 18.93 & & 4.61 & 7.04 & 5.57 & 7.08 \\
    f-VAEGAN-D2 \cite{42} \,{\scriptsize \textbf{CVPR'19}} & 12.77 & 0.95 & 1.77 & 1.91 & & 17.29 & 8.47 & 11.37 & 11.11 & & 4.36 & 2.14 & 2.87 & 2.40 \\
    APN \cite{41} \,{\scriptsize \textbf{IJCV'22}} & 7.48 & 3.88 & 5.11 & 4.49 & & 28.46 & 16.16 & 20.61 & 16.44 & & 9.84 & 5.76 & 7.27 & 6.34 \\
    CJME \cite{3} \,{\scriptsize \textbf{WACV'20}} & 11.96 & 5.41 & 7.45 & 6.84 & & 48.18 & 17.68 & 25.87 & 20.46 & & 16.06 & 9.13 & 11.64 & 9.92 \\
    AVGZSLNet \cite{4} \,{\scriptsize \textbf{WACV'21}} & 13.02 & 2.88 & 4.71 & 5.44 & & 56.26 & 34.37 & 42.67 & 35.66 & & 14.81 & 11.11 & 12.70 & 12.39 \\
    AVCA \cite{1} \,{\scriptsize \textbf{CVPR'22}} & 32.47 & 6.81 & 11.26 & 8.16 & & 34.90 & 38.67 & 36.69 & 38.67 & & 24.04 & 19.88 & 21.76 & 20.88 \\
    TCAF \cite{5} \,{\scriptsize \textbf{ECCV'22}} & 12.63 & 6.72 & 8.77 & 7.41 & & 67.14 & 40.83 & 50.78 & 44.64 & & 30.12 & 7.65 & 12.20 & 7.96 \\
    AVFS \cite{6} \,{\scriptsize \textbf{IJCNN'23}} & 15.62 & 6.00 & 8.67 & 7.31 & & 54.57 & 36.94 & 44.06 & 41.55 & & 14.41 & 8.91 & 11.01 & 9.15 \\
    Hyper-multiple \cite{7} \,{\scriptsize \textbf{ICCV'23}} & 21.99 & 8.12 & 11.87 & 8.47 & & 43.52 & 39.77 & 41.56 & 40.28 & & 20.52 & \underline{21.30} & 20.90 & \underline{22.18} \\
    KDA \cite{9} \,{\scriptsize \textbf{ARXIV'23}} & 13.30 & 7.74 & 9.78 & 8.32 & & 75.88 & 42.97 & 54.84 & \underline{52.66} & & 37.55 & 10.25 & 17.95 & 11.85 \\
    STFT\cite{43} \,{\scriptsize \textbf{TIP'24}} & 11.74 & 8.83 & 10.08 & 8.79 & & 61.42 & 43.81 & 51.14 & 49.74 & & 25.12 & 9.83 & 14.13 & 9.46 \\
    ClipClap-GZSL \cite{2} \,{\scriptsize \textbf{CVPR'24}} & 27.59 & \underline{11.00} & \underline{15.73} & \underline{11.28} & & 78.20 & 42.57 & 55.13 & 46.01 & & \textbf{44.89} & 19.42 & \underline{27.11} & 21.76 \\
    EZ-GZSL \cite{12} \,{\scriptsize \textbf{ECCV'24}} & 17.26 & 8.68 & 11.55 & 8.92 & & 78.32 & \underline{46.35} & \underline{58.24} & 48.27 & & 38.52 & 12.85 & 19.27 & 13.01 \\
    MSTR \cite{14} \,{\scriptsize \textbf{Neurocomputing'25}} & 13.70 & 5.48 & 7.83 & 6.38 & & \textbf{86.32} & 19.97 & 32.43 & 23.57 & & 22.92 & 9.28 & 13.21 & 9.65 \\
    MDST++ \cite{44} \,{\scriptsize \textbf{TCSVT'25}} & 12.21 & 8.38 & 9.94 & 10.32 & & 65.32 & 44.79 & 53.14 & \textbf{55.22} & & 27.53 & 10.72 & 15.43 & 13.12 \\
    SACMA \cite{45} \,{\scriptsize \textbf{Inf.Fusion'25}} & 13.42 & 6.00 & 8.29 & 6.46 & & 60.15 & 27.46 & 37.71 & 29.07 & & 17.12 & 7.32 & 10.25 & 7.52 \\
    
    \noalign{\global\setlength{\oldaboverulesep}{\aboverulesep}}
    \noalign{\global\setlength{\oldbelowrulesep}{\belowrulesep}}
    \noalign{\global\setlength{\aboverulesep}{0pt}}
    \noalign{\global\setlength{\belowrulesep}{1.5pt}}
    \midrule[1pt]
    \rowcolor{gray!25}
    \textbf{KA-GZSL} & 32.47 & \textbf{11.28} & \textbf{16.74} & \textbf{12.16} & & \underline{83.79} & \textbf{46.68} & \textbf{59.96} & 48.98 & & \underline{42.80} & \textbf{23.18} & \textbf{30.07} & \textbf{24.24}\\
    \bottomrule[2pt]
    \noalign{\global\setlength{\aboverulesep}{\oldaboverulesep}}
    \noalign{\global\setlength{\belowrulesep}{\oldbelowrulesep}}
    
  \end{tabular}
\end{table*}


\subsubsection{Comparison methods}
We compare the proposed method with 5 classic ZSL methods and 13 audio–visual GZSL methods. The ZSL methods include DeViSE\cite{38}, ALE\cite{39}, SJE\cite{40}, APN\cite{41}, and f-VAEGAN-D2\cite{42}. Among them, DeViSE, ALE, SJE, and APN are discriminative zero-shot learning methods based on label embeddings and are typically used in image tasks to learn the compatibility between visual features and semantic embeddings. f-VAEGAN-D2 is a generative ZSL method that mitigates the missing sample problem by conditionally generating unseen class features. To adapt these models to the audio–visual task, the aforementioned ZSL baselines take both audio and visual features as inputs in our experiments rather than relying solely on a single visual feature.


The audio–visual GZSL comparison methods include CJME\cite{3}, AVGZSLNet\cite{4}, AVCA\cite{1}, TCAF\cite{5}, AVFS\cite{6}, Hyper-multiple\cite{7}, KDA\cite{9}, STFT\cite{43}, ClipClap-GZSL\cite{2}, EZ-GZSL\cite{12}, MSTR\cite{14}, MDST++\cite{44}, and SACMA\cite{45}. CJME and AVGZSLNet align modalities in a shared embedding space, with AVGZSLNet additionally using cross-modal decoding and reconstruction. AVCA and TCAF model cross-modal temporal interactions, whereas AVFS synthesizes unseen features with contrastive learning and Hyper-multiple uses hyperbolic geometry for hierarchical alignment. KDA combines LLM-generated descriptions with distribution alignment. STFT, MSTR, and MDST++ combine SNNs and transformers to model spatiotemporal and sparse-motion information; ClipClap-GZSL and EZ-GZSL use pretrained multimodal models with audio–visual–text contrastive alignment. SACMA captures intra- and intermodal relations through self- and cross-attention with compositional contrastive and cosine-similarity losses.


\begin{table*}[!t]
  \centering
  \caption{Influence of KAS, AVCE, and AKG in KA-GZSL.}
  \label{tab:component_ablation}
  \begin{tabular}{l cccc c cccc c cccc}
    \toprule[2pt]
    \multirow{2}{*}{\textbf{Model}}
    & \multicolumn{4}{c}{\text{VGGSound-GZSL}$^{\textit{cls}}$} & & \multicolumn{4}{c}{\text{UCF-GZSL}$^{\textit{cls}}$} & & \multicolumn{4}{c}{\text{ActivityNet-GZSL}$^{\textit{cls}}$} \\
    \cmidrule(lr){2-5} \cmidrule(lr){7-10} \cmidrule(lr){12-15}
    & S & U & HM & ZSL & & S & U & HM & ZSL & & S & U & HM & ZSL \\
    \midrule[1pt]
    W/o \textbf{K}nowledge-\textbf{A}ugmented \textbf{S}emantics & 28.67 & 9.39 & 14.15 & 8.66 & & 73.82 & 38.24 & 50.38 & \textbf{52.93} & & 38.13 & 19.79 & 26.06 & 22.42 \\
    W/o \textbf{A}udio--\textbf{V}isual \textbf{C}ontext \textbf{E}ncoder & 18.11 & 7.25 & 10.35 & 7.94 & & 50.64 & 16.37 & 24.74 & 20.94 & & 19.65 & 11.26 & 14.32 & 15.49 \\
    W/o \textbf{A}ction \textbf{K}nowledge \textbf{G}raph & 15.27 & 6.15 & 8.77 & 8.53 & & 45.36 & 16.53 & 24.23 & 15.38 & & 23.65 & 18.37 & 20.68 & 12.12 \\
    \noalign{\global\setlength{\oldaboverulesep}{\aboverulesep}}
    \noalign{\global\setlength{\oldbelowrulesep}{\belowrulesep}}
    \noalign{\global\setlength{\aboverulesep}{0pt}}
    \noalign{\global\setlength{\belowrulesep}{1.5pt}}
    \midrule[1pt]
    \rowcolor{gray!25}
    \textbf{KA-GZSL} & \textbf{32.47} & \textbf{11.28} & \textbf{16.74} & \textbf{12.16} & & \textbf{83.79} & \textbf{46.68} & \textbf{59.96} & 48.98 & & \textbf{42.80} & \textbf{23.18} & \textbf{30.07} & \textbf{24.24} \\
    \bottomrule[2pt]
    \noalign{\global\setlength{\aboverulesep}{\oldaboverulesep}}
    \noalign{\global\setlength{\belowrulesep}{\oldbelowrulesep}}
  \end{tabular}
\end{table*}

\begin{table*}[!t]
  \centering
  \caption{Influence of different prompts in KAS.}
  \label{tab:prompt}
  \begin{tabular}{l cccc c cccc c cccc}
    \toprule[2pt]
    \multirow{2}{*}{\textbf{Prompt}} 
    & \multicolumn{4}{c}{\text{VGGSound-GZSL}$^{\textit{cls}}$} & & \multicolumn{4}{c}{\text{UCF-GZSL}$^{\textit{cls}}$} & & \multicolumn{4}{c}{\text{ActivityNet-GZSL}$^{\textit{cls}}$} \\
    \cmidrule(lr){2-5} \cmidrule(lr){7-10} \cmidrule(lr){12-15}
    & S & U & HM & ZSL & & S & U & HM & ZSL & & S & U & HM & ZSL \\
    \midrule[1pt]
    None & 11.35 & 5.66 & 7.55 & 6.98 & & 50.26 & 35.83 & 41.84 & 42.16 & & 28.73 & 16.22 & 20.73 & 19.51 \\
    a photo of a person doing \{class label\} & 19.20 & 9.98 & 13.13 & 9.11 & & 79.88 & 44.86 & 57.45 & \textbf{49.23} & & \textbf{45.77} & 20.26 & 28.09 & 22.55 \\
    
    \noalign{\global\setlength{\oldaboverulesep}{\aboverulesep}}
    \noalign{\global\setlength{\oldbelowrulesep}{\belowrulesep}}
    \noalign{\global\setlength{\aboverulesep}{0pt}}
    \noalign{\global\setlength{\belowrulesep}{1.5pt}}
    \midrule[1pt]
    \rowcolor{gray!25}
    \textbf{a photo of a... +“class explanation”} & \textbf{32.47} & \textbf{11.28} & \textbf{16.74} & \textbf{12.16} & & \textbf{83.79} & \textbf{46.68} & \textbf{59.96} & 48.98 & & 42.80 & \textbf{23.18} & \textbf{30.07} & \textbf{24.24} \\
    \bottomrule[2pt]
    \noalign{\global\setlength{\aboverulesep}{\oldaboverulesep}}
    \noalign{\global\setlength{\belowrulesep}{\oldbelowrulesep}}
  \end{tabular}
\end{table*}

\subsection{Quantitative analysis}
Table \ref{tab:sota} compares KA-GZSL with prior methods on the three benchmarks. KA-GZSL obtains the best HM on \text{VGGSound-GZSL}$^{\textit{cls}}$, \text{UCF-GZSL}$^{\textit{cls}}$, and \text{ActivityNet-GZSL}$^{\textit{cls}}$ (16.74\%, 59.96\%, and 30.07\%), compared with 15.73\%, 55.13\%, and 27.11\% for ClipClap-GZSL. For ZSL, KA-GZSL leads on \text{VGGSound-GZSL}$^{\textit{cls}}$ and \text{ActivityNet-GZSL}$^{\textit{cls}}$ with 12.16\% and 24.24\%, exceeding ClipClap-GZSL by 0.88 percentage points and Hyper-multiple by 2.06 percentage points, respectively; on \text{UCF-GZSL}$^{\textit{cls}}$, it reaches 48.98\%, below MDST++ (55.22\%). Overall, the results show improved seen--unseen balance across all three datasets, while the highest absolute HM is obtained on the smaller and more focused \text{UCF-GZSL}$^{\textit{cls}}$ benchmark.


\subsection{Ablation analysis}
\subsubsection{Influence of KAS, AVCE, and AKG in KA-GZSL}
Table \ref{tab:component_ablation} shows that removing KAS, AVCE, or AKG reduces S, U, and HM on all three datasets. Removing AVCE decreases HM by 6.39\%, 35.22\%, and 15.75\% on \text{VGGSound-GZSL}$^{\textit{cls}}$, \text{UCF-GZSL}$^{\textit{cls}}$, and \text{ActivityNet-GZSL}$^{\textit{cls}}$, while removing AKG yields corresponding decreases of 7.97\%, 35.73\%, and 9.39\%. Removing KAS decreases HM by 2.59\%, 9.58\%, and 4.01\%, although it raises \text{UCF-GZSL}$^{\textit{cls}}$ ZSL from 48.98\% to 52.93\%. The complete model nevertheless achieves the highest HM on every dataset, confirming the complementary roles of label enrichment, contextual fusion, and structured knowledge transfer.

\subsubsection{Influence of different prompts in KAS}
As shown in Table \ref{tab:prompt}, when the “class explanation” is integrated into the base template “a photo of a person doing {class label}”, the HM and ZSL scores of the model generally increase across all the datasets. Specifically, on the diverse and complex \text{VGGSound-GZSL}$^{\textit{cls}}$ dataset, the introduction of class explanations significantly increases the HM metric from 13.13\% to 16.74\%; on the \text{ActivityNet-GZSL}$^{\textit{cls}}$ dataset, the HM metric also increases from 28.09\% to 30.07\%. The analysis of different prompt construction strategies on the three benchmark datasets reveals that introducing detailed class explanations plays a significant role in enhancing model performance.

\begin{table*}[!t]
  \centering
  \caption{Influence of different modalities}
  \label{tab:modality}
  \begin{tabular}{c cccc c cccc c cccc}
    \toprule[2pt]
    \multirow{2}{*}{\textbf{Modality}} 
    & \multicolumn{4}{c}{\textbf{\text{VGGSound-GZSL}}$^{\textit{cls}}$} & & \multicolumn{4}{c}{\textbf{\text{UCF-GZSL}}$^{\textit{cls}}$} & & \multicolumn{4}{c}{\textbf{\text{ActivityNet-GZSL}}$^{\textit{cls}}$} \\
    \cmidrule(lr){2-5} \cmidrule(lr){7-10} \cmidrule(lr){12-15}
    & S & U & HM & ZSL & & S & U & HM & ZSL & & S & U & HM & ZSL \\
    \midrule[1pt]
    Audio & 19.26 & 10.28 & 13.41 & 9.86 & & 44.17 & 42.37 & 43.25 & 46.27 & & 12.36 & 8.33 & 9.95 & 7.67 \\
    Visual & 17.23 & 7.56 & 10.51 & 7.66 & & 59.66 & 45.16 & 51.41 & 47.39 & & 37.14 & \textbf{24.83} & 29.76 & 23.94 \\
    
    \noalign{\global\setlength{\oldaboverulesep}{\aboverulesep}}
    \noalign{\global\setlength{\oldbelowrulesep}{\belowrulesep}}
    \noalign{\global\setlength{\aboverulesep}{0pt}}
    \noalign{\global\setlength{\belowrulesep}{1.5pt}}
    \midrule[1pt]
    \rowcolor{gray!25}
    \textbf{Both} & \textbf{32.47} & \textbf{11.28} & \textbf{16.74} & \textbf{12.16} & & \textbf{83.79} & \textbf{46.68} & \textbf{59.96} & \textbf{48.98} & & \textbf{42.80} & 23.18 & \textbf{30.07} & \textbf{24.24} \\
    \bottomrule[2pt]
    \noalign{\global\setlength{\aboverulesep}{\oldaboverulesep}}
    \noalign{\global\setlength{\belowrulesep}{\oldbelowrulesep}}
  \end{tabular}
\end{table*}

% \begin{table*}[htbp]  % table* = 跨双栏，彻底解决多列溢出问题
%   \centering
%   \caption{Influence of weighting coefficients $\alpha$, $\beta$, and $\varphi$ in the loss function.}
%   \label{tab:loss_coeff}
%   % 核心：@{\hspace{2em}} 在三个数据集之间强制增加间距，实现“分开”效果
%   \begin{tabular}{ccc *{4}{c} @{\hspace{2em}} *{4}{c} @{\hspace{2em}} *{4}{c}}
%     \toprule[1pt]  % 顶部粗线，还原原图样式
%     \multirow{2}{*}{Factor} & & & \multicolumn{4}{c}{\text{VGGSound-GZSL}$^{\textit{cls}}$} & \multicolumn{4}{c}{\text{UCF-GZSL}$^{\textit{cls}}$} & \multicolumn{4}{c}{\text{ActivityNet-GZSL}$^{\textit{cls}}$} \\
%     \cmidrule{4-7} \cmidrule{8-11} \cmidrule{12-15}  % 分割三个数据集的表头
%     $\alpha$ & $\beta$ & $\varphi$ & S & U & HM & ZSL & S & U & HM & ZSL & S & U & HM & ZSL \\
%     \midrule[0.8pt]  % 表头与数据之间的中粗线
%     $\times$ & $\times$ & $\times$ & 23.53 & 5.68 & 9.15 & 6.94 & 56.38 & 30.45 & 39.54 & 25.68 & 27.04 & 12.73 & 17.31 & 15.64 \\
%     $\times$ & $\checkmark$ & $\checkmark$ & 24.72 & 6.86 & 10.74 & 7.55 & 69.57 & 39.72 & 50.57 & 38.98 & 31.29 & 18.94 & 23.60 & 17.44 \\
%     $\checkmark$ & $\times$ & $\checkmark$ & 28.55 & 7.27 & 11.59 & 9.13 & 74.94 & \textbf{49.92} & 59.92 & 42.64 & 35.71 & 17.59 & 23.57 & 20.43 \\
%     $\checkmark$ & $\checkmark$ & $\times$ & 28.93 & 6.38 & 10.45 & 8.97 & 68.45 & 42.98 & 52.80 & 41.75 & 40.85 & 20.67 & 27.45 & \textbf{24.76} \\
%     $\checkmark$ & $\checkmark$ & $\checkmark$ & \textbf{32.47} & \textbf{11.28} & \textbf{16.74} & \textbf{12.16} & \textbf{83.79} & \textbf{46.68} & \textbf{59.96} & \textbf{48.98} & \textbf{42.80} & \textbf{23.18} & \textbf{30.07} & \textbf{24.24} \\
%     \bottomrule[1pt]  % 底部粗线，还原原图样式
%   \end{tabular}
% \end{table*}
% -------------------- 表1：损失系数分析 --------------------

\subsubsection{Influence of different modalities}
Table \ref{tab:modality} compares audio-only, visual-only, and joint inputs. Audio-only is stronger on \text{VGGSound-GZSL}$^{\textit{cls}}$ (HM/ZSL 13.41\%/9.86\% versus 10.51\%/7.66\% for visual-only), whereas visual-only is stronger on \text{UCF-GZSL}$^{\textit{cls}}$ and \text{ActivityNet-GZSL}$^{\textit{cls}}$; on \text{ActivityNet-GZSL}$^{\textit{cls}}$, its HM is 29.76\% versus 9.95\% for audio-only. Joint input achieves the highest HM on every dataset, including 59.96\% on \text{UCF-GZSL}$^{\textit{cls}}$, 8.55 percentage points above visual-only. This confirms that audio and visual cues are complementary despite their dataset-dependent strengths.

\subsubsection{Additional ablations in the supplementary material}
We further evaluate the contributions of the three training losses, the text encoders used for knowledge-graph construction, and the weighting coefficients in Eq. \eqref{eq:total_loss}. As reported in the \href{run:Supplementary_Materials.pdf}{supplementary material}, the complete loss yields the highest HM on all three datasets. Removing $L_n$ reduces \text{UCF-GZSL}$^{\textit{cls}}$ HM/ZSL from 59.96\%/48.98\% to 17.38\%/18.14\% and \text{VGGSound-GZSL}$^{\textit{cls}}$ HM from 16.74\% to 7.00\%, while removing $L_p$ or $L_c$ lowers \text{ActivityNet-GZSL}$^{\textit{cls}}$ HM from 30.07\% to 14.32\% or 20.26\%, respectively. Joint CLIP--CLAP encoding achieves HM scores of 16.74\%, 59.96\%, and 30.07\% on the three datasets, outperforming either encoder alone. Jointly using $\alpha$, $\beta$, and $\varphi$ also yields the highest HM on every dataset; removing them individually lowers HM by 9.39 points on \text{UCF-GZSL}$^{\textit{cls}}$, 6.50 points on \text{ActivityNet-GZSL}$^{\textit{cls}}$, and 6.29 points on \text{VGGSound-GZSL}$^{\textit{cls}}$, respectively.


\subsection{Qualitative analysis}
The t-SNE\cite{46} visualization is provided in \href{run:Supplementary_Materials.pdf}{Fig.~3 of the supplementary material}. Across all three datasets, the raw audio and visual features are relatively dispersed, whereas the learned joint audio--visual embeddings form more compact clusters and contain their corresponding class-semantic embeddings. This qualitative result supports the alignment between the learned multimodal representation and class semantics; complete visualizations and discussion are provided in the supplementary material.


\section{Conclusion}
In this manuscript, we propose KA-GZSL to address semantic ambiguity, weak audio--visual correlations, and cross-modal alignment in AVGZSL. First, an LLM expands class labels into fine-grained semantic descriptions to improve unseen-class representation. AVCE then integrates modality-specific features with shared temporal information, while AKG constructs global action and audio knowledge graphs and transfers prior knowledge to audio--visual temporal features through GCN-based cross-graph propagation. Triplet, projection, and reconstruction losses jointly align the multimodal representations with textual semantics. Experiments on \text{VGGSound-GZSL}$^{\textit{cls}}$, \text{UCF-GZSL}$^{\textit{cls}}$, and \text{ActivityNet-GZSL}$^{\textit{cls}}$ demonstrate the effectiveness of KA-GZSL in transferring knowledge to unseen classes. Future work will explore lightweight and robust knowledge-enhancement and cross-modal interaction mechanisms and extend KA-GZSL to more complex real-world multimodal scenarios.


\bibliographystyle{IEEEtran}
\bibliography{sec/references}
\begin{IEEEbiography}[{\includegraphics[width=1in,height=1.25in,clip,keepaspectratio]{figure/wh.jpg}}]{Hu Wang}
entered Chengdu University of Information Technology in 2019 and obtained a Bachelor of Engineering degree from Chengdu University of Information Technology in 2023.  In August 2024, he enrolled in the State Key Laboratory of Public Big Data, a jointly constructed laboratory by the province and the ministry, at Guizhou University for postgraduate studies.  His current research direction focuses on zero-shot Learning, audio–visual Learning.\end{IEEEbiography}
\vspace{-3em}

\begin{IEEEbiography}[{\includegraphics[width=1in,height=1.25in,clip,keepaspectratio]{figure/yj.jpg}}]{Jing Yang}
(Member, IEEE) received the PhD degree in mechanical and electronic engineering from Guizhou University, in 2020. From August 2018 to September 2019, he was awarded a scholarship by the China Scholarship Council (CSC) under the State Scholarship Fund to pursue his study with Oklahoma State University, as a Joint Ph.D. Student with the school of Computer Science and Technology, where he joined the Guoliang Fan’s Group. From October 2022 to October 2023, he was a visiting scholar studying in the team of Professor Guo Minyi (IEEE Fellow) from Shanghai Jiao Tong University. He is currently an assistant professor with the State Key Laboratory of Public Big Data, Guizhou University, China. 
His research interests include high performance computing, task scheduling in various architectures and Open domain visual learning. He has published more than 50 peer-reviewed papers in the related area, including well-archived international journals such as IEEE Transactions on Circuits and Systems for Video Technology, IEEE Transactions on Green Communications and Networking, IEEE Transactions on Intelligent transportation systems, and IEEE Internet of Things Journal. He serves more than 20 prestigious international journals and conferences, such as IEEE\ TC, IEEE\ TMC, IEEE\ TI, IEEE\ TIM, IEEE\ TSC, IEEE\ TSM, IEEE\ TAI, IEEE\ TASE, IEEE\ TOJ.\end{IEEEbiography}
% \vspace{-4em}

\begin{IEEEbiography}[{\includegraphics[width=1in,height=1.25in,clip,keepaspectratio]{figure/rxl.png}}]{Xiaoli Ruan}
 received the PhD degree in information and communication engineering from Yunnan University, in 2021. Participating in the "Jiaxiu Zhiguang" Visiting Scholar Program at Zhejiang University from September 2024 to June 2025. She is currently a lecturer with the State Key Laboratory of Public Big Data, Guizhou University, China. Her research interests include Artificial Intelligence and Multimodal Big Data Fusion. She has published more than 20 peer-reviewed papers in the related area, including well-archived international journals such as IEEE Transactions on Geoscience and Remote Sensing, Knowledge-Based Systems, and IEEE Internet of Things Journal. 
\end{IEEEbiography}

\begin{IEEEbiography}[{
\includegraphics[width=1in,height=1.25in,clip,keepaspectratio]{figure/cyl.jpg}}]{Yuling Chen}
received her BS from Taishan University, Taian, PR China, in 2006; MS from Guizhou
University, Guiyang, PR China, in 2009; Ph.D.
from Guizhou University, Guiyang, PR China, in
2021. She is now a professor at the State Key
Laboratory of Public Big Data, Guizhou University,
Guiyang, PR China. Her recent research interests
include federated learning, blockchain, cryptography
and information safety. received the M.E degrees in
the College of Computer Science and Technology
from Guizhou University, China, in 2018 and 2021. Her research interests are cryptography, blockchain, information security and artifcial intelligence security.\end{IEEEbiography}

\begin{IEEEbiography}[{\includegraphics[width=1in,height=1.25in,clip,keepaspectratio]{figure/yxr.jpg}}]{Xinru Yi}
was admitted to Guizhou University in September 2023 to pursue a bachelor's degree in Mechanical Design, Manufacturing and Automation. She has been awarded the National Endeavor Scholarship, the National Second Prize in the China Robot and Artificial Intelligence Competition, and the Provincial Excellence Award in the Challenge Cup China National College Students' Entrepreneurship Competition. Her research interests include automatic control and machine learning.\end{IEEEbiography}

\begin{IEEEbiography}[{\includegraphics[width=1in,height=1.25in,clip,keepaspectratio]{figure/lcj.jpg}}]{Chengjiang Li}
 received the PhD degree in Condensed Matter Physics from the Department of Physics, University of Oxford, UK, in 2012. He obtained his Bachelor's degree in Microelectronics from the School of Physical Science and Technology, Sichuan University, China, in 2007. His research interests include Quantum Information, New Energy Materials, and Optoelectronic Devices. He has published more than 60 peer-reviewed papers indexed by SCI/EI in high-impact international journals, including Journal of Materials Chemistry A, Light: Science \& Applications, Science China Materials, Advanced Optical Materials, Optics Express, Applied Physics Letters, Journal of Applied Physics, and Nano Letters.
\end{IEEEbiography}

\begin{IEEEbiography}[{\includegraphics[width=1in,height=1.25in,clip,keepaspectratio]{figure/gmy.jpg}}]{Minyi Guo}
(Fellow, IEEE) received the BSc and ME
 degrees in computer science from Nanjing University,
 China, and the PhD degree in computer science from
 the University of Tsukuba, Japan. Additionally, he is
 the Department of Computer Science and Engineering director of the Embedded and Pervasive Computing Center, Shanghai Jiao Tong University. Before
 joining SJTU,he had been a professor with the School
 of Computer Science and Engineering, University of
 Aizu, Japan. He received the national science fund for distinguished young scholars from NSFC in 2007. His present research interests include parallel/distributed computing, compiler optimizations, Big Data, and cloud computing. He has more than 400 publications in major journals and international conferences in these areas. He
 received 7 best/highlight paper awards from international conferences including
 ALSPOS 2017 and ICCD 2018. He is now editor-in-chief of IEEE Transactions
 on Sustainable Computing, and on the editorial board of IEEE Transactions
 on Parallel and Distributed Systems, IEEE Transactions on Cloud Computing,
 and Journal of Parallel and Distributed Computing. He is the Zhiyuan chair
 professor, and an ACM distinguished member.\end{IEEEbiography}

 
\end{document}
